% ==========================================================================
% ch05.tex — dihasilkan otomatis dari website/chapters/ch05.qmd
% Boleh disunting manual. 'sync' hanya menimpa bila .qmd berubah DAN .tex
% belum disunting; kalau keduanya berubah, dilewati (pakai --force).
% qmd-sha1: fb804da995a7
% tex-sha1: 846582b624be
% ==========================================================================
% >>>>> ISI OTOMATIS (jangan hapus baris ini) >>>>>
\chapter{Missing Values \& Outliers (Reproducible, Berbasis Pipeline)}\label{missing-values-outliers-reproducible-berbasis-pipeline}

Ketiadaan data dan kehadiran \emph{outlier} merupakan kenyataan empiris yang senantiasa dijumpai saat bekerja dengan data mentah di dunia nyata. Algoritma pemelajaran mesin tidak dapat memproses nilai kosong pada saat kalkulasi gradien atau metrik jarak, sedangkan \emph{outlier} ekstrem yang menyimpang jauh dari pusat distribusi dapat mendistorsi estimasi parameter model. Sayangnya, tindakan pembersihan data baris per baris di luar alur sistematis sering kali mengubah variansi asli, menyembunyikan anomali yang valid, dan memicu \emph{data leakage} jika parameter pengisi dihitung sebelum pemisahan data (\emph{data split}) dilakukan.

Bab ini mendeskripsikan fondasi penanganan nilai kosong dan \emph{outlier} secara sistematis, \emph{reproducible}, dan terintegrasi di dalam objek \emph{pipeline}. Pembahasan mencakup mekanisme hilangnya data (MCAR, MAR, MNAR) sebagai asupan diagnostik sebelum otomatisasi, pemilihan imputasi sederhana melawan imputasi iteratif multivariat berbasis model, pemanfaatan status kekosongan sebagai \emph{feature} representatif (\emph{missing indicator}), serta transformasi \emph{outlier} menggunakan parameter statistik yang tangguh (\emph{robust}). Seluruh proses ini dikemas ke dalam alur kerja terstruktur untuk menjamin integritas distribusi data dari tahap pelatihan hingga fase inferensi di lingkungan produksi.

\section{\texorpdfstring{\emph{Missing Values} \& \emph{Outlier} (Reproducible, Berbasis \emph{pipeline})}{Missing Values \& Outlier (Reproducible, Berbasis pipeline)}}\label{missing-values-outlier-reproducible-berbasis-pipeline}

\textbf{Mekanisme \emph{Missingness}: MCAR, MAR, dan MNAR.}\label{mekanisme-missingness-mcar-mar-dan-mnar}

Sebagian besar algoritma \emph{machine learning} tidak dapat memproses nilai kosong secara langsung, sehingga data yang tidak lengkap harus disesuaikan sebelum tahap pemodelan. Sebelum menyusun strategi imputasi di dalam \emph{pipeline}, alasan di balik hilangnya data tersebut perlu diinspeksi. Ketiadaan nilai sering kali merepresentasikan kejadian faktual yang membawa informasi spesifik tentang prosedur perekaman observasi. Tahap diagnostik ini sangat penting sebelum menerapkan transformasi otomatis.

Secara matematis, mekanisme hilangnya data (\emph{missingness mechanism}) didefinisikan melalui probabilitas bersyarat dari matriks indikator kekosongan. Misalkan terdapat sekumpulan data \(Y\) yang terdiri dari bagian teramati \(Y_{obs}\) dan bagian kosong \(Y_{miss}\). Jika kita mendefinisikan matriks indikator \(R\) dengan nilai 1 berarti data hilang dan 0 berarti teramati, mekanisme \emph{missingness} dinyatakan sebagai fungsi probabilitas \(P(R \mid Y_{obs}, Y_{miss}, \phi)\). Dalam fungsi tersebut, \(R\) merepresentasikan matriks indikator ketiadaan data, \(Y_{obs}\) adalah nilai atribut yang berhasil terekam, \(Y_{miss}\) adalah nilai atribut yang kosong, dan \(\phi\) merupakan parameter tidak diketahui yang menentukan proses hilangnya data. Berdasarkan hubungan ketergantungan ini, statistika membagi hilangnya data ke dalam tiga kategori utama.

Kategori pertama adalah \emph{Missing Completely At Random} (MCAR). Pada kasus ini, kekosongan terjadi murni secara kebetulan dan peluang hilangnya data tidak dipengaruhi oleh variabel lain. Bentuk distribusinya disederhanakan menjadi \(P(R \mid Y_{obs}, Y_{miss}, \phi) = P(R \mid \phi)\). Contohnya, sensor pengukur angin gagal mengirimkan data sesaat akibat pemadaman listrik acak. Pada skenario MCAR, porsi data yang utuh tetap menjadi sampel representatif dari populasi awal.

Kategori kedua adalah \emph{Missing At Random} (MAR). Peluang terjadinya nilai kosong memiliki probabilitas bersyarat dengan variabel lain di dalam \emph{dataset}, tetapi tidak dipengaruhi oleh nilai atribut yang hilang itu sendiri. Persamaannya diukur menjadi \(P(R \mid Y_{obs}, Y_{miss}, \phi) = P(R \mid Y_{obs}, \phi)\). Sebagai contoh, responden survei berusia lanjut sering mengosongkan riwayat aktivitas fisik. Pola kekosongan pada riwayat olahraga tersebut pada dasarnya berkorelasi matematis dengan rentang usia responden. Memisahkan pengamatan berdasarkan umur akan menghasilkan probabilitas kekosongan yang independen pada setiap kelompok batas usianya.

Kategori ketiga adalah \emph{Missing Not At Random} (MNAR). Probabilitas hilangnya suatu nilai terikat langsung dengan isi dari nilai yang tidak teramati tersebut (\(Y_{miss}\)). Ketergantungan ini membuat persamaan probabilitas tidak dapat disederhanakan. Kasus klasiknya adalah survei pendapatan, di mana responden dengan angka kekayaan ekstrem sering kali menolak melapor demi privasi. Absennya data merupakan dampak langsung dari nominal kekayaan itu sendiri, sehingga pada kondisi MNAR, lokasi kekosongan justru mengandung sinyal pemodelan yang kuat.

\begin{figure}[htbp]\centering
\includegraphics[width=\linewidth,height=0.72\textheight,keepaspectratio]{ch05-fig-1}
\caption{Perbandingan Pola Hilangnya Data pada Mekanisme MCAR, MAR, dan MNAR}
\label{fig:ch05-fig-1}
\end{figure}

Membedakan ketiga profil mekanisme di atas mencegah injeksi bias pada fase pelatihan. Jika kita mengisi data yang tergolong MNAR menggunakan nilai rata-rata, kita akan mendistorsi batas distribusi asli dan menyamarkan keberadaan kelompok ekstrem dalam sampel.

Penerapan fungsi pengamatan statistik memfasilitasi proses inspeksi ini. Pengujian berbasis validasi observasional ini dapat dicapai melalui visualisasi matriks kekosongan yang direpresentasikan oleh fungsi pustaka \passthrough{\lstinline!missingno!}. Korelasi kosong (\emph{nullity}) komparatif di antara dua lajur parameter menandakan kehadiran relasi struktural bertipe MAR maupun MNAR. Pada evaluasi formal, uji statistik MCAR Little dioperasikan untuk meninjau secara mekanis rasio asumsi deviasi \emph{missingness}. Meskipun indikasi matematis memberikan acuan awal, pendekatan kuantitatif ini harus diverifikasi berdampingan dengan evaluasi konteks domain faktual data terkait. Jika kondisi absen bertipe MNAR teridentifikasi, parameter kosong tersebut berfungsi menjadi sinyal observasi mandiri. Beberapa arsitektur \emph{deep learning} di desain khusus memetakan nilai kosong (\emph{missing value}) sebagai bagian indikator mekanisme pemelajaran internal untuk menyelaraskan akurasi prediksi asimetris.

Operasi imputasi \emph{missing values} memerlukan pendekatan yang sistematis. Pemahaman konseptual, inspeksi visual, serta pengetahuan konteks domain merupakan prasyarat mutlak sebelum mengubah teknik penanganan nilai kosong menjadi komponen \emph{pipeline} permanen.

\section{\texorpdfstring{Imputasi Sebagai Komponen \emph{Pipeline}}{Imputasi Sebagai Komponen Pipeline}}\label{imputasi-sebagai-komponen-pipeline}

Praktik yang sering keliru dalam pengembangan model adalah memposisikan penanganan data kosong sebagai tahap pembersihan di awal proses. Data diisi menggunakan rata-rata atau median dari keseluruhan \emph{dataset}, baru kemudian dibagi menjadi \emph{split training} dan \emph{split test}. Pendekatan ini memicu \emph{data leakage} karena tahap imputasi tidak terisolasi dari data uji. Imputasi adalah transformasi struktural yang harus tertanam di dalam \emph{pipeline} pemodelan.

Fenomena \emph{data leakage} terjadi karena perhitungan statistik merekam informasi dari kelompok evaluasi. Jika sekumpulan data \(X\) merupakan gabungan data latih dan data uji, persamaan rata-rata imputasi yang terdistorsi adalah:

\[ \mu_{bocor} = \frac{1}{|X_{train}| + |X_{test}|} \left( \sum_{x \in X_{train}} x + \sum_{x \in X_{test}} x \right) \]

Suku \(\sum_{x \in X_{test}} x\) memastikan bahwa parameter pengisi pada data latih dihitung dengan menyertakan komponen distribusi uji. Akibatnya, model menyesuaikan fungsi internalnya berdasarkan rentang angka observasi masa depan yang direpresentasikan saat sesi inisialisasi pelatihan. Hasilnya, metrik kalkulasi skor tampak sangat akurat pada data eksperimental, namun akan mengalami tingkat kesalahan klasifikasi yang signifikan saat diaplikasikan pada baris \emph{dataset} di operasional lingkungan \emph{production}.

\begin{figure}[htbp]\centering
\includegraphics[width=\linewidth,height=0.72\textheight,keepaspectratio]{ch05-fig-2}
\caption{Alur Imputasi Data Utuh yang Menyebabkan Kebocoran}
\label{fig:ch05-fig-2}
\end{figure}

\begin{figure}[htbp]\centering
\includegraphics[width=\linewidth,height=0.72\textheight,keepaspectratio]{ch05-fig-3}
\caption{Alur Imputasi Aman Menggunakan Pipeline}
\label{fig:ch05-fig-3}
\end{figure}

Untuk menghindari \emph{data leakage}, modul imputasi (seperti \passthrough{\lstinline!SimpleImputer!} atau \passthrough{\lstinline!KNNImputer!} pada pustaka Scikit-Learn) harus diperlakukan sama seperti model prediktif, melalui pemisahan fase \emph{fit} dan \emph{transform}. Nilai pengisi dihitung eksklusif dari \emph{split training}:

\[ \mu_{benar} = \frac{1}{|X_{train}|} \sum_{x \in X_{train}} x \]

Siklus fungsi arsitektur dari parameter indikasi imputasi ini memaksa evaluasi statistik untuk dikerjakan independen. Pertama, kerangka kerja terisolasi berimplikasi pada fungsi komputasi \emph{cross-validation}. Parameter nilai estimasi diformulasikan ulang pada blok data terpisah pelatihan (\emph{fold}), menghindarkan tumpukan lipatan kelompok uji (\emph{validation fold}) dari kebocoran sinyal referensi perhitungan parameter aslinya. Kedua, kerangka operasional imputasi yang konsisten meniadakan kendala komputasi pada fungsional algoritma di tingkat operasional *production*. Varians observasi tabel asing yang didatangkan untuk dievaluasi oleh parameter metode prediktif secara fungsional dinilai hanya merujuk acuan distribusi pelatihan yang sebelumnya disematkan oleh modul imputer bawaan. Ketiga, prosedur penyisipan yang memusat meredam distorsi komputasi untuk mekanisme multivariat interaktif. Algoritma seperti prediksi klasifikasi antar kolom \emph{features} diharuskan berjalan independen memanfaatkan murni \emph{split training} observasinya.

Meskipun beberapa kajian literatur menyimpulkan bahwa imputasi median tak-diawasi (*unsupervised*) yang tidak beroperasi melalui isolasi \emph{cross-validation} sesekali memperkecil perhitungan tingkat ukuran varians dengan membebankan penalti fungsional parameter metrik observasi pada ukuran parameter bias (Jaeger dkk., 2020), standar mutlak prapemrosesan harus dijaga. Penggabungan logika algoritma imputer dikerjakan konsisten hanya sesudah fungsi isolasi prapemrosesan rentang pemisahan *split* parameter observasi diterapkan sebagai komponen terpadu variabel *transformer* pada urutan parameter komputasi \emph{pipeline}.

\section{Imputasi Sederhana vs.~Berbasis Model}\label{imputasi-sederhana-vs.-berbasis-model}

Pemilihan algoritma pengisian nilai kosong bergantung pada mekanisme \emph{missingness} (MCAR, MAR, atau MNAR). Pendekatan imputasi terbagi ke dalam dua metode: metode univariat dasar yang memprioritaskan kecepatan kalkulasi dan metode multivariat yang berusaha mempertahankan struktur distribusi data asli.

\textbf{Imputasi Univariat (Sederhana).}\label{imputasi-univariat-sederhana}

Imputasi sederhana mengisi sel data kosong menggunakan metrik statistik tunggal dari kolom yang sama, seperti rata-rata (\emph{mean}), median, atau modus.

Karakteristik metode pendekatan imputasi komputasi univariat didasarkan murni pada tingkat kemudahan komputasi algoritma pemrosesannya. Kalkulasi matematis bekerja dengan singkat dengan menyajikan evaluasi fungsi linear sekuensial. Tetapi, kelebihan tersebut hadir mengabaikan fungsi analisis silang (\emph{cross-correlation}) karena setiap variabel kolom beroperasi mutlak secara independen. Misalnya jika rentang dimensi parameter berat badan dinormalisasi berbasis indeks median global distribusi dari *dataset*, kalkulasi estimasi angka yang hilang tersebut tidak dipandu rujukan data variabel fungsional kelompok umur maupun proporsi dimensi ukuran tinggi secara spesifik dari sampel baris pengamatan itu. Imputasi angka titik berpusat juga berdampak signifikan dengan menekan simpangan angka (\emph{variance shrinkage}) data tersebut. Konsentrasi jumlah titik yang sama pada garis median secara proporsional menekan pelebaran sebaran data populasi awal. Rekayasa parameter nilai varians mendistorsi rasio perbandingan galat dimensi matematis (\emph{standard error}) ketika \emph{machine learning} melacak konvergensi prediksi komputasi titik sasarannya.

\textbf{Imputasi Multivariat (Berbasis Model).}\label{imputasi-multivariat-berbasis-model}

Untuk mempertahankan korelasi struktural antar-\emph{feature}, imputasi multivariat memperlakukan kolom berstatus \emph{missing} sebagai variabel target prediksi, sementara lajur kolom utuh difungsikan secara bergantian sebagai matriks penjelas.

Dua algoritma klasik utama di kelas multivariat ini adalah K-Nearest Neighbors (KNN) dan Imputasi Iteratif (MICE). Algoritma KNN mengisi sel renggang melalui rata-rata nilai dari gabungan tetangga terdekatnya. Proses pencariannya mengeksekusi parameter jarak spesifik hanya dari rentang variabel yang bebas sel kosong. Di sisi lain, \emph{Multiple Imputation by Chained Equations} (MICE) menyusun rantai prediksi dengan melatih fungsi regresi pada tiap-tiap target kolom secara berurutan dalam skema \emph{round-robin}.

\begin{figure}[htbp]\centering
\includegraphics[width=\linewidth,height=0.72\textheight,keepaspectratio]{ch05-fig-4}
\caption{Siklus Prediksi Berantai pada Algoritma MICE}
\label{fig:ch05-fig-4}
\end{figure}

Secara matematis, prediksi fungsi target pada skema imputasi iteratif dikalkulasikan dengan formula \(\hat{x}_{i,j}^{(t)} = f_j \left( x_{i,1}^{(t)}, \dots, x_{i,j-1}^{(t)}, x_{i,j+1}^{(t-1)}, \dots, x_{i,p}^{(t-1)} \right)\). Dalam persamaan tersebut, \(\hat{x}_{i,j}^{(t)}\) menyatakan angka proyeksi pada sel baris ke-\(i\) di kolom \(j\) untuk siklus ke-\(t\). Variabel \(f_j\) mewakili algoritma pemelajaran mesin tersupervisi (contohnya \emph{Random Forest} atau regresi) yang ditugaskan khusus membedah varians kolom \(j\). Model ini mencerna deret data secara kontinu: meminjam pembaruan nilai kelompok indeks \(1\) sampai \(j-1\) dari iterasi terkini \((t)\), lalu memanfaatkan sisa rentang dari \(j+1\) hingga dimensi puncak \(p\) dari siklus masa lalu \((t-1)\). Modul arsitektur pustaka modern seperti \passthrough{\lstinline!IterativeImputer!} umumnya dikonfigurasi untuk memberikan prediksi nilai tunggal (\emph{single imputation}) demi mereduksi kerumitan \emph{pipeline} komputasi inferensi, walau hal ini menghilangkan visibilitas atas ketidakpastian angka penambalan (\emph{imputation uncertainty}) yang biasanya dikelola lewat metodologi MICE murni.

\textbf{Pendekatan Generatif Berbasis Deep Learning.}\label{pendekatan-generatif-berbasis-deep-learning}

Pada saat matriks \emph{feature} berkembang menjadi dimensi raksasa yang mengandung distribusi dependensi tak-linear, perhitungan \emph{deep learning} menghasilkan variasi penyelesaian komputasi imputasi fungsional tersupervisi. Sebagai gambaran struktural, model jaringan \emph{Masked Autoencoder} mengubah parameter kekosongan pengamatan menjadi segmen lapisan (\emph{mask}) untuk menstimulasi fungsi pengenalan mandiri \emph{self-supervised} dalam menalar sebaran rasio varians asli pada data berdimensi utuh. Pada model *diffusion* probabilistik (\emph{Diffusion Models}), pergeseran matematis metode penyebaran transmisi *noise* diadaptasi untuk menghitung kemungkinan pengisian nilai \emph{feature} dengan pendekatan estimasi yang menyerupai fungsi *baseline* rasio asli data. Modifikasi desain untuk jaringan deteksi pola MNAR (seperti fungsional model \passthrough{\lstinline!not-MIWAE!}) melepaskan belenggu probabilitas matematis dari sistem acak (MAR) pada skema evaluasi *training* konvensional. Model arsitektur ini memetakan observasi parameter \emph{missingness} sebagai mekanisme terintegrasi ke dalam rancangan pembobotan prediksi matriks observasi fungsional multi-layer.

Dominasi tingkat performa arsitektur jaringan generasi baru ini hampir selalu menuntut biaya komputasi yang tinggi. Banyak pakar komparatif metrik modern, direpresentasikan dalam dokumen analisis Näf (2026), merekomendasikan komputasi fungsional imputasi stokastik berbantuan regresi pohon berjejer (\emph{MICE-Random Forest}) lantaran margin selisih kapabilitas prediksinya kerap kali sangat tipis dibandingkan model generatif yang kompleks. Keberhasilan parameter probabilitas klasik menjustifikasi fungsinya baik sebagai \emph{baseline} fungsional maupun parameter penyeimbang pada optimisasi utilitas ruang arsitektur generatif modern.

\section{\texorpdfstring{\emph{Missing Indicator} Sebagai Sinyal}{Missing Indicator Sebagai Sinyal}}\label{missing-indicator-sebagai-sinyal}

Ketika berhadapan dengan data yang tidak lengkap, proses imputasi secara otomatis menghapus sebuah informasi struktural: fakta bahwa observasi tersebut awalnya kosong. Ketidakhadiran sebuah nilai sering kali menyimpan sinyal prediktif yang lebih kuat daripada nilai tebakan yang dimasukkan oleh algoritma imputasi.

Untuk mempertahankan informasi ini, kita mengonstruksi \emph{feature} biner tambahan yang disebut \emph{missing indicator}. \emph{Feature} ini memisahkan operasi pengisian nilai dari status ketersediaan data awal. Secara matematis, untuk sebuah \emph{feature} \(x^{(j)}\), indikator \(m^{(j)}\) didefinisikan melalui persamaan berikut:

\[
m_i^{(j)} = \begin{cases}
1, & \text{jika } x_i^{(j)} \text{ bernilai NaN (kosong)} \\
0, & \text{jika } x_i^{(j)} \text{ terobservasi}
\end{cases}
\]

Dalam rumusan ini, variabel \(m_i^{(j)}\) mewakili elemen indikator biner untuk observasi ke-\(i\) pada \emph{feature} ke-\(j\), sedangkan \(x_i^{(j)}\) merepresentasikan nilai asli dari data sebelum melewati tahapan imputasi.

\begin{figure}[htbp]\centering
\includegraphics[width=\linewidth,height=0.72\textheight,keepaspectratio]{ch05-fig-5}
\caption{Alur Transformasi Satu Fitur Asli Menjadi Fitur Berimputasi dan Missing Indicator}
\label{fig:ch05-fig-5}
\end{figure}

Pendekatan ini sangat relevan saat menangani mekanisme data \emph{Missing Not At Random} (MNAR), di mana pola kekosongan memiliki korelasi dengan variabel target. Sebagai contoh, pada model penilaian risiko kredit, ketiadaan data ``nomor telepon tempat kerja'' sering kali bukan akibat kegagalan pencatatan sistem. Kekosongan tersebut dapat mengindikasikan riwayat pekerjaan yang tidak stabil, yang secara empiris merupakan sinyal probabilitas gagal bayar yang tinggi.

Keberadaan \emph{missing indicator} mengubah cara model membaca \emph{dataset} yang telah diimputasi. Algoritma akan menggunakan \emph{feature} hasil imputasi untuk mengidentifikasi pola pada populasi umum, sembari di saat yang sama memanfaatkan matriks indikator biner ini untuk menyesuaikan probabilitas prediksi berdasarkan fenomena ketiadaan informasi tersebut. Pada model berbasis pohon (seperti \emph{Random Forest} atau \emph{Gradient Boosting}), tersedianya indikator ini memungkinkan algoritma membuat partisi percabangan logika (\emph{split}) spesifik yang membedakan observasi dengan nilai asli dari observasi dengan nilai tebakan hasil perkiraan.

Secara teknis, komputasi indikator ini umumnya dirangkai dalam satu paket dengan objek imputasi di dalam \emph{pipeline} (seperti mengaktifkan parameter \passthrough{\lstinline!add_indicator=True!} pada \passthrough{\lstinline!SimpleImputer!} di pustaka Scikit-Learn). Pembentukan indikator tunduk pada batasan fiksasi \emph{fit-transform} untuk mencegah \emph{data leakage}. Pertama, \emph{feature} indikator ini dikunci secara fiksasi saat pelatihan; matriks biner hanya diregistrasi untuk kolom-kolom yang secara empiris terdeteksi memiliki nilai kosong di dalam set pelatihan (\emph{training set}). Kedua, terdapat kekakuan saat inferensi; jika sebuah kolom terisi penuh di data pelatihan namun secara mengejutkan muncul dengan nilai kosong pada data baru di lingkungan produksi, \emph{pipeline} tidak akan mengizinkan pembentukan kolom indikator baru secara dadakan, melainkan hanya akan menjalankan fungsi imputasinya saja.

Penggunaan \emph{missing indicator} memiliki konsekuensi komputasi yang menuntut evaluasi struktural terukur. Praktisi wajib menggunakannya secara cermat dengan mengantisipasi dua konsekuensi spasial. Risiko pertama berupa penambahan indikator redundan yang mengekspansi skala dimensi fitur secara berlebih (karena imputasi MCAR meratakan proporsi tanpa menambah acuan korelasi spesifik) tanpa menambahkan tingkat sinyal signifikansi terhadap variabel \emph{dataset}. Risiko sekunder terletak pada pembentukan multikolinearitas struktural fungsional matriks linear. Kondisi \emph{dataset} yang menghapus data indikator dalam kombinasi parameter seragam berbarengan secara independen (\emph{co-missing patterns})--misalnya tanggal masuk kontrak tercatat bersama kolom keluarnya dari rentang operasional--berakibat parameter indikator yang sama pada fungsi korelasinya mengalami garis keselarasan konstan absolut. Kelemahan kalkulasi multikolinearitas mengganggu kestabilan arah konvergensi optimisasi saat dioperasikan memakai metode linear regresi atau algoritma parametrik murni lainnya.

\section{\texorpdfstring{\emph{Outlier}: Deteksi vs.~Penanganan}{Outlier: Deteksi vs.~Penanganan}}\label{outlier-deteksi-vs.-penanganan}

Selain penanganan \emph{missing values}, nilai ekstrem atau \emph{outlier} merupakan tantangan utama dalam merepresentasikan distribusi populasi. \emph{Outlier} kerap diperlakukan sebagai data rusak yang harus segera dihapus. Praktik pembersihan reaktif ini berisiko menghilangkan varians data. Rekayasa \emph{feature} yang sistematis membagi manajemen \emph{outlier} ke dalam dua tahap operasional terpisah: deteksi dan penanganan.

\textbf{Deteksi: Mengobservasi Anomali.}\label{deteksi-mengobservasi-anomali}

Tahap deteksi murni bertujuan menetapkan batas kenormalan suatu \emph{feature}. Algoritma hanya melabeli atau memberikan skor pada nilai yang secara matematis berada di luar distribusi mayoritas, tanpa melakukan modifikasi nilai.

Ekosistem analitik data mengelompokkan pendekatan deteksi \emph{outlier} ke dalam beberapa tingkatan perhitungan algoritma komputasi. Pendekatan pertama didasari pengukuran univariat, yang menganalisis rasio pusat distribusi melalui penerapan batasan ambang matematis interkuartil (IQR). Pada observasi ruang multi-dimensi fungsional, metode pendekatan kedua yaitu analisis multivariat dieksekusi dengan mendayagunakan metode \emph{Isolation Forest} dan perhitungan jarak batas ruang variabel \emph{Local Outlier Factor} (LOF). Evaluasi \emph{Isolation Forest} membedah rentang observasi menjadi jarak dimensi pemisahan titik pencilan dari sampel mayoritas dalam kelompok *dataset*, sementara LOF menentukan skor jarak spasial menggunakan nilai proksimal kepadatan dimensi dari matriks bertingkat tersebut. Untuk kebutuhan komputasi skala rasio tanpa perhitungan metrik absolut manual, tingkat yang ketiga memanfaatkan pendekatan probabilitas observasional. Implementasi \emph{Empirical-CDF-based Outlier Detection} (ECOD) menghasilkan batas deteksi anomali pada ruang kumulatif titik dimensi spasial yang menekan kebutuhan parameter penyetelan fungsional fungsi distribusi secara independen.

Implementasi antarmuka dalam pustaka seperti Scikit-Learn juga membedakan skenario analisis anomali ke dalam dua metode. \emph{Outlier detection} melabeli data pelatihan yang memang sudah terkontaminasi anomali sejak awal (umumnya menggunakan fungsi \passthrough{\lstinline!fit_predict!}). Di sisi lain, \emph{novelty detection} membangun profil normal secara eksklusif dari data bersih, lalu menggunakannya untuk mendeteksi penyimpangan hanya pada observasi baru (misalnya dengan memicu parameter \passthrough{\lstinline!novelty=True!} pada LOF).

Untuk data univariat, rata-rata dan deviasi standar mudah terdistorsi oleh nilai ekstrem itu sendiri. Metode yang lebih tangguh (\emph{robust}) mengganti deviasi standar dengan skor Z yang dimodifikasi (\emph{modified z-score}) yang diturunkan dari Deviasi Absolut Median (MAD):

\[ MAD = \text{median}(|x_i - \text{median}(X)|) \]

\[ M_i = \frac{0.6745 \cdot (x_i - \text{median}(X))}{MAD} \]

Dalam persamaan ini, variabel \(x_i\) merepresentasikan nilai observasi \emph{feature} ke-\(i\), sedangkan \(\text{median}(X)\) merupakan nilai tengah dari seluruh observasi \emph{feature} \(X\). Variabel \(MAD\) mengukur variabilitas dari penyebaran data di sekitar median yang kebal terhadap tarikan \emph{outlier}. Terakhir, \(M_i\) adalah metrik anomali observasi; konstanta \(0.6745\) digunakan untuk menyelaraskan metrik ini agar sepadan dengan estimasi deviasi standar pada distribusi normal. Secara empiris, observasi dengan nilai \(|M_i| > 3.5\) umumnya dilabeli sebagai anomali.

\begin{figure}[htbp]\centering
\includegraphics[width=\linewidth,height=0.72\textheight,keepaspectratio]{ch05-fig-6}
\caption{Perbandingan Alur Eksekusi Data Cleaning Statis vs Transformasi Dinamis}
\label{fig:ch05-fig-6}
\end{figure}

\textbf{Penanganan: Konteks Domain.}\label{penanganan-konteks-domain}

Setelah modul matematika mengkalkulasi skor ekstrem, keputusan penyelesaian anomali tidak dipandu murni lewat rentang distribusi komputasi. Konteks riil operasional (*domain knowledge*) dari rentang *dataset* secara spesifik harus mendikte bentuk konversi dan substitusinya. Pada skenario kesalahan sistem pencatatan pengukuran instrumen, seperti observasi umur konsumen negatif atau angka ukuran sensor pada kisaran mustahil secara biologi manusia, baris tersebut secara efektif dikalkulasikan sebagai fitur kolom bernilai \emph{missing value}. Solusinya bisa direalisasikan ke fungsi pemotongan titik maksimum batas operasional (\emph{capping}) dari algoritma. Kasus anomali yang kedua terjadi karena pengamatan tersebut adalah pola sinyal asli matriks data. Perekaman jumlah transfer atau ukuran komputasi batas jutaan transaksi bisa berarti pergerakan nasabah *corporate* khusus (yang secara sah terindikasi *fraud*) pada evaluasi deteksi klasifikasi bank. Kesalahan fatal pada parameter modifikasi imputasi yang tidak seimbang untuk mendepak representasi \emph{outlier} tanpa dasar akan mendistorsi prediksi evaluasi akhir dan parameter optimisasi batas fungsi \emph{machine learning}. Pemisahan titik isolasi fungsional parameter batas pengurungan dari metode intervensi perhitungan nilai menjamin kalkulasi parameter asli dan mengembalikan signifikansi rentang *machine learning* secara utuh.

\section{\texorpdfstring{Transformasi \emph{Robust} untuk \emph{Outlier}}{Transformasi Robust untuk Outlier}}\label{transformasi-robust-untuk-outlier}

\emph{Outlier} yang faktual dan valid menimbulkan dilema. Menghapus observasi membuang informasi dan mengurangi ukuran sampel data. Sebaliknya, membiarkan nilai ekstrem apa adanya dapat mendistorsi pelatihan. Algoritma yang sensitif terhadap magnitudo (seperti regresi linier atau metode berbasis jarak) rentan terhadap efek ini karena fungsi kerugiannya memberikan penalti besar pada titik ekstrem. Akibatnya, garis regresi tertarik kuat ke arah \emph{outlier} hanya untuk meminimalkan kerugian pada titik tersebut, sehingga mengorbankan akurasi mayoritas data. Transformasi \emph{robust} meredam pengaruh nilai ekstrem ini tanpa menghilangkan informasi posisinya.

Terdapat tiga metode utama untuk menangani nilai ekstrem melalui pendekatan transformasi.

Metode pertama adalah pemangkasan (\emph{capping} atau \emph{winsorizing}). Daripada membuang baris data, kita menetapkan ambang batas atas dan batas bawah; setiap nilai yang melampauinya dipangkas dan disamakan dengan batas tersebut. Sebagai contoh, jika batas ditetapkan pada persentil ke-99, observasi pendapatan yang melebihi titik tersebut diubah menjadi sama persis dengan nilai persentil itu. Individu terkait tetap berada di kelompok teratas, tetapi magnitudo absolutnya terkendali. Implementasi yang sejalan dengan praktik \emph{pipeline} yang benar akan mempelajari batas pemangkasan ini murni dari data latih dan menerapkannya kembali secara persis pada saat pemrosesan data uji.

Ambang batas pemangkasan ini umumnya dihitung menggunakan tiga jenis pendekatan matematis. Pendekatan berbasis kuantil memotong distribusi menurut proporsi absolut, misalnya menyisihkan 1\% data terendah dan tertinggi. Pendekatan berbasis Rentang Antarkuartil (IQR) menggunakan rentang antar-persentil sebagai basis turunan pagar anomali Tukey. Sementara itu, pendekatan berbasis Deviasi Absolut Median (MAD) menjadi alternatif \emph{robust} dari standar deviasi konvensional. Standar deviasi biasa akan terdistorsi jika dihadapkan pada \emph{outlier} ekstrem karena nilai ekstrem tersebut otomatis ikut membesarkan hasil kalkulasi rata-rata dan varians; perhitungan berbasis MAD kebal terhadap distorsi tersebut karena ia bersandar kuat pada pengukuran median.

Persamaan \emph{modified z-score} menggunakan rumusan:

\[ M_i = \frac{0.6745 \cdot (x_i - \tilde{x})}{\text{MAD}} \]

Dalam formulasi tersebut, variabel \(M_i\) mewakili skor modifikasi untuk observasi ke-\(i\), \(x_i\) adalah nilai titik data, \(\tilde{x}\) adalah median dari distribusi \emph{feature}, dan \(\text{MAD}\) merupakan deviasi absolut median itu sendiri (yakni median dari selisih absolut setiap observasi terhadap \(\tilde{x}\)). Faktor \(0.6745\) digunakan semata-mata untuk menyejajarkan skor dengan rentang distribusi normal baku. Titik data yang melewati rentang absolut \(|M_i| > 3.5\) biasanya dipangkas atau ditandai sebagai \emph{outlier}.

\begin{figure}[htbp]\centering
\includegraphics[width=\linewidth,height=0.72\textheight,keepaspectratio]{ch05-fig-7}
\caption{Arsitektur Pipeline Pemrosesan Data: Aliran X\_train dan X\_test}
\label{fig:ch05-fig-7}
\end{figure}

Selain pemangkasan, penyesuaian distribusi \emph{feature} melalui penskalaan juga lazim digunakan. Metode kedua adalah penskalaan \emph{robust} (\emph{robust scaling}). Penskalaan standar berbasis \emph{z-score} menjadi sangat rentan karena rata-rata dan standar deviasinya terseret oleh \emph{outlier}. Penskalaan \emph{robust} memusatkan data dengan median dan menskalakannya menggunakan IQR. Pengukuran ini murni memakai titik data di pusat distribusi (persentil ke-25 sampai ke-75), yang membuat skala \emph{feature} tetap stabil meskipun terdapat deretan nilai ekstrem di ujung rentang distribusi.

Metode ketiga adalah pemampatan distribusi (\emph{log / power transform}). Memampatkan skala terbukti efektif memfasilitasi data yang didominasi nilai kecil namun memiliki ekor panjang ke sisi kanan. Transformasi logaritmik atau penerapan \emph{power transform} menekan nilai-nilai dominan ekstrem agar rentang magnitudonya bergeser mendekati sisa distribusi mayoritas. Operasi ini mengecilkan jarak absolut antar-observasi tanpa mendisrupsi urutan ordinal; nilai maksimal tetap di puncak, sehingga model tidak terdistorsi akibat besaran asimetris saat pelatihan.

\section{\texorpdfstring{Batas Antara \emph{Data Cleaning} dan Rekayasa \emph{Feature}}{Batas Antara Data Cleaning dan Rekayasa Feature}}\label{batas-antara-data-cleaning-dan-rekayasa-feature}

Dalam alur kerja ilmu data, batasan antara \emph{data cleaning} dan rekayasa \emph{feature} kerap beririsan karena keduanya bertugas mengubah wujud data mentah sebelum dieksekusi oleh model \emph{machine learning}. Pemahaman tentang batas operasional kedua proses ini berguna untuk merancang \emph{pipeline} yang kebal terhadap \emph{data leakage}, sekaligus mendefinisikan ruang lingkup teknis yang dibahas dalam buku ini.

\emph{Data cleaning} berfokus pada koreksi fakta untuk memastikan sebuah \emph{dataset} merepresentasikan keadaan dunia nyata secara akurat. Validitas dari pembersihan data bersifat mutlak dan tidak akan berubah meskipun kita mengganti algoritma pemodelan. Pertama, perbaikan ini terfokus pada fakta yang salah, format yang tidak konsisten, atau duplikasi entitas. Kedua, titik eksekusinya bersifat statis dan global; ia dieksekusi satu kali pada keseluruhan data sebelum proses pemodelan atau pembagian data dimulai. Contoh tindakannya meliputi mengurai teks \emph{string} ``1,200.50'' agar terbaca secara benar sebagai nilai numerik, menghapus observasi usia 150 tahun yang mustahil secara biologis, atau melakukan deduplikasi terhadap baris data yang tercatat berulang akibat kesalahan sistem pencatatan.

Di sisi lain, rekayasa \emph{feature} beroperasi di ranah representasi. Fokus utamanya bukan lagi menguji kebenaran fakta dasar, melainkan memanipulasi bentuk data agar algoritma dapat mengekstrak sinyal darinya dengan efisien. Pertama, manipulasi ini bertujuan mengubah distribusi, skala, atau format representasi sebagai keputusan desain yang disesuaikan secara khusus dengan kebutuhan algoritma tertentu. Kedua, titik eksekusinya bersifat dinamis dan harus diikat kuat di dalam \emph{pipeline} pembelajaran untuk menjaga batasan isolasi data latih. Contoh tindakannya mencakup mengaplikasikan \emph{robust scaler} pada sebuah \emph{feature} agar metrik jarak SVM dapat konvergen, mencetak \emph{missing indicator} baru untuk menampung sinyal informasi, atau melakukan imputasi nilai yang hilang berdasarkan pola distribusi observasi latih.

Kriteria pembatas operasional paling tegas antara kedua proses ini bersandar pada literatur kebocoran data (IBM, 2024), yakni dengan mengajukan satu pengujian: apakah transformasi tersebut bergantung pada pembagian \emph{train/test split}?

Secara matematis, perbedaan sifat eksekusi ini dapat dirumuskan melalui definisi fungsi pemrosesan:

\[
X_{bersih} = C(X_{mentah}) \quad \text{vs} \quad X_{\text{feature}} = T_{\theta}(X_{bersih})
\]

Fungsi pembersihan \(C\) diterapkan secara seragam pada seluruh observasi \(X_{mentah}\) tanpa mengamati interaksi antar-baris. Sebaliknya, transformasi \emph{feature} \(T\) dikendalikan oleh matriks parameter \(\theta\) (misalnya batas parameter kuartil pada \emph{scaler} atau nilai rata-rata pada imputer) yang secara eksklusif dipelajari hanya dari subset data latih. Menghitung \(\theta\) menggunakan data pengujian adalah asal mula terjadinya \emph{data leakage}.

Batas operasional ini juga membatasi cakupan materi yang dieksplorasi di dalam bab-bab berikutnya. Buku ini tidak membahas tutorial manipulasi teks mentah, penggunaan fungsi \emph{regex} untuk menyortir format kotor, atau instruksi teknik deduplikasi basis data. Titik berat kita dikhususkan sepenuhnya pada desain dan penerapan operasional \(T_{\theta}\), yakni rekayasa representasi yang berinteraksi langsung dengan mekanisme pembelajaran mesin.

\section{\texorpdfstring{Studi Kasus: \emph{Pipeline} Imputasi dan Penanganan \emph{Robust} yang \emph{Reproducible}}{Studi Kasus: Pipeline Imputasi dan Penanganan Robust yang Reproducible}}\label{studi-kasus-pipeline-imputasi-dan-penanganan-robust-yang-reproducible}

Penerapan taktik penanganan \emph{missing values} dan \emph{outlier} secara terpisah pada tahap prapemrosesan rentan memicu \emph{data leakage}. Jika kalkulasi batas \emph{outlier} atau rata-rata imputasi dilakukan sebelum membagi set pelatihan dan pengujian, informasi dari set pengujian akan ikut merembes dan memengaruhi model. Untuk mencegah kontaminasi ini, seluruh komponen transformasi harus disatukan ke dalam arsitektur \emph{pipeline}.

Sebagai contoh, tinjau sebuah \emph{dataset} medis untuk prediksi risiko penyakit yang memiliki dua anomali struktural:

\begin{enumerate}
\def\labelenumi{\arabic{enumi}.}
\tightlist
\item
  \textbf{Ketiadaan data yang tidak acak:} \emph{Feature} ``jumlah tes darah'' mengandung \emph{missing values} bertipe \emph{Missing Not At Random} (MNAR). Pasien berisiko tinggi lebih sering menjalani observasi laboratorium. Kekosongan data di kolom ini merepresentasikan ketiadaan tes karena pasien relatif sehat.
\item
  \textbf{Ekor distribusi yang panjang:} \emph{Feature} ``kadar protein'' memiliki nilai ekstrem yang valid secara biologis, bukan akibat kesalahan sensor. Jika tidak dibatasi, magnitudo ekstrem ini akan mengganggu konvergensi dan stabilitas skala model linier.
\end{enumerate}

Berdasarkan diagnosis di atas, kita merancang spesifikasi \emph{pipeline} berurutan:

\begin{enumerate}
\def\labelenumi{\arabic{enumi}.}
\tightlist
\item
  \textbf{Ekstraksi indikator:} Menambahkan \passthrough{\lstinline!MissingIndicator!} (atau mengaktifkan argumen \passthrough{\lstinline!add_indicator=True!} pada objek imputer) untuk menyalin lokasi sel yang kosong menjadi sebuah \emph{feature} biner baru.
\item
  \textbf{Imputasi berbasis model:} Menjalankan algoritma \passthrough{\lstinline!IterativeImputer!} atau \passthrough{\lstinline!KNNImputer!} untuk menaksir nilai numerik pada sel yang kosong dengan memetakan korelasi antar-atribut medis.
\item
  \textbf{Pembatasan ekstrem (\emph{capping}):} Menerapkan objek \passthrough{\lstinline!Winsorizer!} untuk memangkas magnitudo nilai di atas batas persentil ke-95.
\item
  \textbf{Penyesuaian skala (\emph{scaling}):} Melakukan standardisasi dengan \passthrough{\lstinline!RobustScaler!} menggunakan perhitungan median dan \emph{Interquartile Range} (IQR), sehingga skala \emph{feature} akhir tetap stabil terhadap sisa residu \emph{outlier}.
\end{enumerate}

Secara matematis, \emph{pipeline} mengeksekusi komposisi fungsi transformasi secara berantai pada data observasi (\(X_{\text{in}}\)). Parameter transformasi di setiap tahap akan dikunci dengan mengandalkan statistik yang dipelajari mutlak dari data latih melalui persamaan \(X_{\text{out}} = T_n(T_{n-1}( \dots T_1(X_{\text{in}} \mid \theta_1) \dots \mid \theta_{n-1}) \mid \theta_n)\). Pada perumusan tersebut, \(T_i\) adalah representasi matematis dari fungsi transformasi pada tahap ke-\(i\) (seperti fungsi pemotongan distribusi atau fungsi pergeseran skala), dan \(\theta_i\) mewakili parameter statistik spesifik (seperti batas kuantil, median, atau bobot imputasi) yang dikalkulasi secara eksklusif saat algoritma memproses \emph{split} pelatihan. Pemanggilan fungsi \passthrough{\lstinline!fit()!} pada blok \emph{pipeline} akan merekam pola indikator variabel, distribusi relasional, dan batas persentil tanpa melihat satu pun observasi dari set pengujian. Saat dieksekusi di dalam skema \emph{cross-validation}, struktur ini memaksa prosedur imputasi dan estimasi statistik beroperasi secara terisolasi murni di dalam setiap porsi lipatan (\emph{fold}).

Sebagai catatan evaluasi teknis, literatur menunjukkan satu pengecualian: metode imputasi univariat statis tanpa pengawasan (seperti imputasi nilai rata-rata biasa) yang dieksekusi sebelum pemisahan lipatan \emph{cross-validation} memang terkadang memunculkan bias yang dapat diabaikan sembari mempercepat komputasi. Namun, untuk metode imputasi iteratif dan tersupervisi, keseluruhan kalkulasinya wajib diletakkan secara ketat di dalam \emph{pipeline} guna menghindari ilusi akurasi evaluasi yang fiktif.

\begin{lstlisting}[language=Python]
import numpy as np
import pandas as pd
from sklearn.base import BaseEstimator, TransformerMixin
from sklearn.experimental import enable_iterative_imputer  # Wajib diimpor sebelum IterativeImputer
from sklearn.impute import IterativeImputer
from sklearn.preprocessing import RobustScaler
from sklearn.pipeline import Pipeline

class Winsorizer(BaseEstimator, TransformerMixin):
    """
    Transformer kustom untuk memotong (cap/clip) nilai outlier ekstrem 
    berdasarkan nilai persentil batas bawah dan batas atas.
    """
    def __init__(self, lower_quantile=0.05, upper_quantile=0.95):
        self.lower_quantile = lower_quantile
        self.upper_quantile = upper_quantile
        self.lower_bounds_ = None
        self.upper_bounds_ = None
        
    def fit(self, X, y=None):
        X_df = pd.DataFrame(X)
        self.lower_bounds_ = X_df.quantile(self.lower_quantile).values
        self.upper_bounds_ = X_df.quantile(self.upper_quantile).values
        return self
        
    def transform(self, X):
        X_df = pd.DataFrame(X).copy()
        for col_idx in range(X_df.shape[1]):
            X_df.iloc[:, col_idx] = np.clip(
                X_df.iloc[:, col_idx], 
                self.lower_bounds_[col_idx], 
                self.upper_bounds_[col_idx]
            )
        return X_df.values

# Membangun arsitektur pipeline pemrosesan data medis terintegrasi
medical_pipeline = Pipeline([
    ('imputer', IterativeImputer(max_iter=10, random_state=42, add_indicator=True)),
    ('outlier_clipper', Winsorizer(lower_quantile=0.01, upper_quantile=0.95)),
    ('scaler', RobustScaler())
])

# Simulasi data medis pasien (misalnya: Tekanan Darah, Kadar Gula Darah, Kadar Protein)
np.random.seed(42)
X_raw = np.random.normal(loc=100, scale=20, size=(100, 3))
X_raw[0, 0] = 400.0  # Sisipkan outlier ekstrem atas
X_raw[1, 1] = 2.0    # Sisipkan outlier ekstrem bawah
X_raw[5:15, 0] = np.nan  # Sisipkan missing values (NaN)
X_raw[12:20, 2] = np.nan

# Melatih pipeline sekaligus mentransformasikan data latih secara aman
X_processed = medical_pipeline.fit_transform(X_raw)

print(f"Dimensi data asli: {X_raw.shape}")
print(f"Dimensi data terproses (termasuk kolom indikator missing): {X_processed.shape}")
print(f"Nilai maksimum setelah pemrosesan: {np.max(X_processed):.4f}")
\end{lstlisting}

\section{Sintesis Keputusan dan Prinsip Akhir}\label{sintesis-keputusan-dan-prinsip-akhir}

Penanganan nilai kosong dan pencilan adalah perpaduan antara kehati-hatian statistik, pemahaman konteks domain, dan kedisiplinan integrasi \emph{pipeline}.

\textbf{Strategi Pengisian Nilai Kosong.} Diagnosis pola kekosongan terlebih dahulu sebelum memilih metode. Pada skenario MCAR, imputasi univariat sederhana (rata-rata/median) cukup memadai dan efisien. Jika korelasi antar-variabel kuat (MAR), terapkan imputasi multivariat (seperti KNN atau MICE) agar struktur distribusi asli tetap terjaga. Pada kondisi MNAR, ketiadaan data sering kali menyimpan informasi penting; pertahankan sinyal tersebut dengan mengonstruksi \emph{feature} biner pendukung berupa \emph{missing indicator}.

\textbf{Manajemen Pencilan.} Deteksi pencilan menggunakan parameter yang kebal terhadap kontaminasi ekstrem (seperti median dan MAD untuk data univariat, atau \emph{Isolation Forest} untuk data berdimensi tinggi). Keputusan intervensi (apakah dipotong lewat \emph{capping/winsorization}, ditransformasi lewat metode \emph{robust}, atau diperlakukan sebagai \emph{missing values}) wajib disesuaikan dengan konteks domain untuk mencegah hilangnya sinyal anomali yang bernilai tinggi (seperti pada deteksi transaksi penipuan).

\textbf{Integrasi Bersih dalam Pipeline.} Seluruh kalkulasi statistik (seperti rata-rata imputasi, median, dan batas kuantil pemotongan) harus diisolasi di dalam objek \emph{pipeline}. Parameter ini dipelajari (\emph{fit}) murni dari himpunan data latih dan diterapkan kembali saat \emph{transform} pada data uji atau data baru di lingkungan produksi. Batasan tegas ini menjamin reproduksibilitas komputasi sekaligus melenyapkan risiko kebocoran data.

Demonstrasi interaktif dan latihan penanganan data renggang secara terstruktur ini dapat diuji secara langsung menggunakan skrip dan data pendamping pada repositori buku ini.
